% =============================================================================
% Method — ACFL journal manuscript
% =============================================================================

\section{Adversarially-Conditioned Federated Learning}
\label{sec:method}

We propose \emph{Adversarially-Conditioned Federated Learning} (ACFL), a server-side defense framework for federated cyber-physical inference on unmanned aircraft systems.
ACFL extends the standard FL training loop at three modular insertion points: pre-aggregation update validation (Layer~1), trust-weighted robust aggregation (Layer~2), and inference-time spoofing detection (Layer~3).
Layers~1 and~2 operate during training; Layer~3 uses the hardened global model at deployment.
This section describes the threat model, FL baseline, and each defense layer; evaluation protocol is given in Section~\ref{sec:experiments}.

% ---------------------------------------------------------------------------
\subsection{Threat Model}
\label{sec:method-threat}

We consider adversaries that compromise one or more FL clients and manipulate \emph{learning-plane} submissions---local training data, outgoing model updates, or both---while the server and benign clients follow the protocol.
The adversary may know the global model architecture and FL hyperparameters (gray-box) but not private ACFL tuning or probe contents.
Attacks are applied every round unless noted otherwise.

Table~\ref{tab:attacks-method} summarizes the learning-plane attacks evaluated in this work and their mapping to ACFL mechanisms.

\begin{table}[t]
\centering
\caption{Learning-plane attacks and ACFL countermeasures.}
\label{tab:attacks-method}
\small
\setlength{\tabcolsep}{3pt}
\begin{tabular}{@{}p{0.22\columnwidth}p{0.40\columnwidth}p{0.30\columnwidth}@{}}
\toprule
Attack & Mechanism & Primary ACFL layer \\
\midrule
Gradient scaling & Amplified update $\mathbf{w}_g + \alpha\Delta$ & L1 subspace + L2 filter \\
Byzantine & Arbitrary/noisy update & L2 coordinate robustness \\
Backdoor & Data poison + trigger update & L1 counterfactual probes \\
Label flip & Corrupted local labels & L2 mission-context trust \\
Model replacement & Anti-direction update & L1 + L2 combined \\
\bottomrule
\end{tabular}
\end{table}

Network-plane attacks (GNSS spoofing, ADS-B injection, MAVLink waypoint manipulation) are out of scope for the quantitative results in Section~\ref{sec:experiments} but motivate Layer~3 and the counterfactual probe scenarios used in Layer~1~\cite{humphreys2008spoofing,strohmeier2011adsb,yaacoub2020uav}.

% ---------------------------------------------------------------------------
\subsection{Federated Learning Baseline}
\label{sec:method-fl}

Let $N$ clients hold private datasets $\mathcal{D}_1,\ldots,\mathcal{D}_N$.
At round $t$, the server broadcasts global weights $\mathbf{w}_g^{(t)}$; client $i$ performs local training to obtain $\mathbf{w}_i^{(t)}$ and uploads an update.
Standard FedAvg aggregates
\begin{equation}
  \mathbf{w}_g^{(t+1)} = \sum_{i \in \mathcal{S}_t} \frac{n_i}{\sum_{j \in \mathcal{S}_t} n_j}\,\mathbf{w}_j^{(t)},
  \label{eq:fedavg}
\end{equation}
where $\mathcal{S}_t$ is the set of participating clients and $n_i$ is the local sample count.
A compromised client may submit $\tilde{\mathbf{w}}_i^{(t)} \neq \mathbf{w}_i^{(t)}$ instead of its honest update.

ACFL replaces Eq.~\eqref{eq:fedavg} with a trust-weighted robust aggregate (Section~\ref{sec:method-l2}) after scoring each submission (Section~\ref{sec:method-l1}).

% ---------------------------------------------------------------------------
\subsection{Model Architecture and Training}
\label{sec:method-model}

All experiments use \texttt{TinyPacketCNN}, a lightweight 1D convolutional classifier designed for onboard packet-window classification under tight compute and radio budgets.
It replaces the drift-fairroad default 2-layer LSTM ($\sim$54k parameters) with a fixed-receptive-field CNN ($\sim$13.7k trainable parameters, $\approx$4$\times$ smaller uploads).

\paragraph{Input representation.}
Each training example is a sliding window of $T{=}50$ consecutive MAVLink packets with $F{=}5$ numeric features per timestep: \texttt{SrcPort}, \texttt{DstPort}, \texttt{Length}, \texttt{MsgID}, and \texttt{Protocol}.
Packets are $z$-score normalized per agent using statistics fit on train-mission rows only (the leading 80\% of each train mission under the temporal split; held-out eval missions reuse the same agent-level normalization).
Windows are extracted with stride 25 (50\% overlap).
A window is labeled \emph{attack} if any packet in the span carries \texttt{Label}${=}1$; otherwise it is benign.
Agents with attack-bearing train missions oversample positive windows (with replacement) to a per-agent target of 25\% attack prevalence before local training.

\paragraph{Architecture.}
The model accepts input $\mathbf{X} \in \mathbb{R}^{B \times T \times F}$ and transposes to channel-first form for \texttt{Conv1d}.
The feature extractor comprises three convolutional stages with batch normalization and ReLU:
(i)~$\mathrm{Conv1d}(F{\rightarrow}24,\,k{=}5)$ followed by $\mathrm{MaxPool1d}(2)$,
(ii)~$\mathrm{Conv1d}(24{\rightarrow}48,\,k{=}3)$,
and (iii)~$\mathrm{Conv1d}(48{\rightarrow}64,\,k{=}3)$,
then $\mathrm{AdaptiveAvgPool1d}(1)$ to collapse the temporal axis.
A dropout layer ($p{=}0.15$) precedes a linear head that emits a single logit for binary attack classification.
Table~\ref{tab:model} lists layer widths and training hyperparameters.

\paragraph{Local and federated training.}
Each FL round, a client trains on 16 windows drawn sequentially from its local buffer, performing 2 local epochs with mini-batch size 16.
Optimization uses Adam ($\eta{=}0.002$) on \texttt{BCEWithLogitsLoss} with per-client class balancing: $\texttt{pos\_weight} = n_{\mathrm{neg}}/n_{\mathrm{pos}}$ computed from the client's current training windows.
No weight decay or gradient clipping is applied; weights use PyTorch default initialization.
The server broadcasts the global state dict each round; clients upload FP16+gzip-compressed weights ($\approx$25\,KB per round on the drone profile).
After 12 federated rounds, checkpoints are evaluated on held-out mission windows.
Predictions use $\sigma(\text{logit})$; F1 is computed at threshold 0.5, with an optional validation-grid threshold search over $\{0.05,\ldots,0.95\}$ when it improves held-out F1.
Layer~3 spoofing detection reuses the same architecture and scoring function on live packet windows (Section~\ref{sec:method-l3}).

% ---------------------------------------------------------------------------
\subsection{Layer 1: Counterfactual Update Validation}
\label{sec:method-l1}

Layer~1 assigns each client $i$ a trust score $T_i^{(1)} \in [T_{\min}, 1]$ before aggregation, where $T_{\min}$ is a configurable floor (0.03 in the drone profile).
Let $\boldsymbol{\delta}_i = \mathrm{vec}(\mathbf{w}_i - \mathbf{w}_g)$ denote the flattened update.
Three complementary signals are combined:

\paragraph{Subspace projection.}
A randomized low-rank basis $\mathbf{U} \in \mathbb{R}^{d \times r}$ ($r{=}16$) is drawn once per round.
Each update is projected and its residual norm compared to the cohort median; large residuals indicate anomalous structure consistent with backdoor or outlier updates.

\paragraph{Temporal consistency.}
Update norms and directions are tracked across rounds per client.
Abrupt jumps relative to the client's own history increase poisoning risk, capturing Byzantine and scaling attacks that destabilize inter-round trajectories.

\paragraph{Counterfactual probes.}
The server holds a fixed probe set $\mathcal{P}$ of input windows spanning benign flight and attack scenarios (GNSS spoofing, ADS-B injection, MAVLink waypoint tampering, jamming, backdoor trigger).
For each candidate update, the server temporarily applies the update to a clone of $\mathbf{w}_g$ and measures operational harm---the drop in benign-scenario confidence and rise in attack-scenario confidence on $\mathcal{P}$.
Harm beyond tolerance $\tau$ is amplified by gain $\gamma$ and mapped to a lower trust score.
The \texttt{backdoor\_trigger} scenario is double-weighted because it directly tests backdoor persistence.
On the drone production profile, each scenario is capped at 16 windows and only attack-relevant scenarios are evaluated (Section~\ref{sec:exp-probe-config}), cutting Layer~1 cost while preserving Byzantine/label-flip F1 and accepting a small backdoor F1 tradeoff.

The raw Layer~1 score is a weighted sum (weights 0.50 / 0.20 / 0.30 for subspace, temporal, counterfactual when probes are available), passed through a soft filter, and smoothed with exponential moving average ($\alpha{=}0.35$) so a single benign-looking round cannot fully rehabilitate a poisoner.

% ---------------------------------------------------------------------------
\subsection{Layer 1/2 Gating}
\label{sec:method-gating}

A key design choice is separating \emph{telemetry} Layer~1 scores from \emph{aggregation} Layer~1 scores when Layer~2 peer-robustness gating is active.
For Byzantine-style attacks, Layer~1 scores that would over-penalize benign clients are replaced at aggregation time by a neutral baseline ($T_{\min}$) while peer-aware norm statistics apply a soft cap derived from coordinate robustness.
Raw Layer~1 scores are still logged for trust AUC evaluation.
The production gate is binary.
An adaptive alternative opens that gate only for a Layer~1 cohort penalty
\begin{equation}
  c_i = \max\!\left(0,\; R_i - T_i^{(1)}\right),
  \qquad
  g_i =
  \begin{cases}
    \mathrm{clip}\!\left(\frac{c_i - \tau}{w},\, 0,\, 1\right)
      & \text{if at least two clients have } c_i > \tau, \\
    0 & \text{otherwise,}
  \end{cases}
  \label{eq:adaptive-gate}
\end{equation}
where \(R_i\) is the peer-robustness score, \(\tau{=}0.25\), and \(w{=}0.25\).
When the cohort gate stays shut, aggregation keeps the norm-capped Layer~1 score,
\(T_i = (1-g_i)\,T_i^{\mathrm{L1,capped}} + g_i\,T_i^{\mathrm{gated}}\).
If a norm hard-reject is withheld because dropping the client would break quorum, and that client still satisfies \(T_i^{(1)} - R_i > \tau\), the round is skipped and the global model is held.
The same hold applies when exactly one participant has \(T_i^{(1)} - R_i > \tau\), including a round whose norm ratio stays under the hard-reject threshold.
A client that Layer~1 has already scored down stays in the aggregate, which is the gradient-scaling and Byzantine case. Two clients with that gap also stay in the aggregate.
Under this adaptive path the single-client guard uses the arbitrated trust, so a lone update with high raw Layer~1 and low robustness is still held.
This path is opt-in via \texttt{--acfl-adaptive-l1l2-gating} and is off in the reported production stack.
For backdoor runs, aggregation uses the production \textbf{min-blend} policy: $\min(T_{\mathrm{neutral\_cap}}, T_{\mathrm{raw\_cap}})$, where $T_{\mathrm{neutral}}$ applies Byzantine-style norm soft-capping and $T_{\mathrm{raw}}$ incorporates full Layer~1 probe scores.
In addition, an attack-scoped \textbf{geometry hard-reject} forms a multi-Krum shortlist of the $k{=}2$ lowest peer-distance scores, then retains the shortlist member with the lowest absolute counterfactual trigger $P(\text{attack})$ (\emph{shortlist-veto}), with Krum score as a within-margin tie-break.
This uses ACFL's probe set to demote geometrically plausible but trigger-hot updates---the failure mode of pure Krum/\emph{keep-best} selection---while preserving Layer~1 telemetry and Layer~2 mission-context weighting on the selected update.
A probe-harm hard-reject remains available but is L1-trust-gated (high-trust clients are not dropped on harm alone), because relative trigger-harm outliers are often benign under this attack.
After a finite Layer~2 aggregation, a \textbf{post-aggregation trigger scrub} (on by default for backdoor) takes a few SGD steps on the global weights using server-held clean probes: stamped trigger windows are driven toward benign ($P(\text{attack}){\to}0$), with benign and real-attack anchors plus an $\ell_2$ ball around the pre-scrub weights to limit clean-task drift.
Section~\ref{sec:exp-scrub-ablation} shows that shortlist-veto alone leaves mean trigger near classical levels ($\approx$0.18), while veto+scrub drops it to $\approx$0.05--$0.06$ with flat clean F1.
For label-flip runs, soft Layer~2 mission trust alone underperforms Krum because flipped labels produce near-normal update norms.
Production ACFL therefore enables an attack-scoped \textbf{shortlist-clean} geometry gate: form a multi-Krum shortlist ($k{=}2$), keep the geometric winner unless another shortlist member's clean-probe quality exceeds it by a margin ($0.025$), then aggregate the selected update under Layer~2 weighting.
We also evaluated a \emph{shortlist-clean-blend} variant that averages the two cleanest shortlist peers; it did not widen the mean F1 gap versus Krum and is disabled by default.
We also evaluated a backdoor-specific \emph{probe-aware soft cap} (trigger-only harm with tighter tolerance); Section~\ref{sec:exp-probe-cap} reports that ablation---it did not outperform min-blend and is disabled in the default configuration.
Tensor aliasing during probe evaluation is avoided by cloning state dictionaries before counterfactual forward passes.

% ---------------------------------------------------------------------------
\subsection{Layer 2: Trust-Weighted Robust Aggregation}
\label{sec:method-l2}

Layer~2 performs two stages.

\paragraph{Stage 1: Coordinate-wise robust pre-filtering.}
For each parameter coordinate, the server computes the median and trimmed mean across client updates.
Clients whose coordinates deviate beyond a dynamic threshold---calibrated from benign-round history and scaled by coordinate deviation threshold (2.5) and norm-ratio thresholds (2.0 general, 1.2 for two-client rounds)---are soft-filtered: their updates are blended toward the robust statistic rather than hard-dropped, preserving effective sample size under bandwidth-constrained avionics links.

\paragraph{Stage 2: Mission-context trust weighting.}
Each client receives a statistical robustness score $R_i$ from Stage~1 and a mission-context score $M_i$ derived from client-reported telemetry headers (GNSS/IMU consistency, MAVLink state consistency, link health, local validation loss, participation history).
These are combined with the gated Layer~1 score:
\begin{equation}
  \tilde{T}_i = \bigl(0.6\, R_i + 0.4\, M_i\bigr)\, T_i^{(1)},
  \qquad
  w_i = \max\!\left(T_{\min},\; \tilde{T}_i^{\,\rho}\right),
  \label{eq:l2-trust}
\end{equation}
where $\rho{=}1.25$ is the trust-weight power that sharpens separation once a weak client is identified.
Final aggregation is $w_i$-weighted over filtered updates, normalized to sum to one.
Optional guards withhold aggregation when participation falls below a minimum quorum or when a sole participant appears suspicious (pairwise inversion guard).

% ---------------------------------------------------------------------------
\subsection{Layer 3: Inference-Time Spoofing Detection}
\label{sec:method-l3}

Layer~3 is not part of the FL aggregation loop.
At inference, the hardened global model scores live packet windows with predicted $P(\text{attack})$, optionally calibrated against Layer~1 benign/attack probe anchors, and fused with optional GNSS/MAVLink consistency headers from client mission context.
Regression (GPS/IMU) models retain cross-modal command-vs-state inconsistency scoring.
If the fused score exceeds a threshold, the system recommends \texttt{safe\_hold}; otherwise \texttt{continue}.
Layer~3 is implemented and exposed via a \texttt{/detect\_spoofing} endpoint; held-out SITL evaluation reports aggregate AUROC $\approx 0.96$ (Section~\ref{sec:exp-limitations}).

% ---------------------------------------------------------------------------
\subsection{Complexity and Deployment}
\label{sec:method-complexity}

Layer~1 probe evaluation dominates server-side cost but runs on a compact model (13{,}745 parameters) with a small probe set.
Weight transfer uses FP16 compression and gzip ($\approx$25\,KB per round on the drone profile).
All defense logic is server-side; clients only add standard telemetry headers, requiring no change to onboard training code beyond the existing drift-fairroad client.
